\documentclass[11pt]{article}
\usepackage[T1]{fontenc}
\usepackage{lmodern,amsmath,amssymb,amsthm,mathtools}
\usepackage[margin=1in]{geometry}
\usepackage{microtype,enumitem,xcolor}
\usepackage[numbers,sort&compress]{natbib}
\usepackage{hyperref}
\hypersetup{colorlinks=true,linkcolor=blue!45!black,citecolor=blue!45!black,
urlcolor=blue!45!black,pdftitle={Small Lyapunov ratios and the structure of Berry--Esseen maximizers},pdfauthor={Logan Bell}}
\setlength{\emergencystretch}{2em}
\numberwithin{equation}{section}
\newtheorem{theorem}{Theorem}[section]
\newtheorem{proposition}[theorem]{Proposition}
\newtheorem{lemma}[theorem]{Lemma}
\newtheorem{corollary}[theorem]{Corollary}
\theoremstyle{remark}\newtheorem{remark}[theorem]{Remark}
\DeclareMathOperator{\Var}{Var}\DeclareMathOperator{\Cov}{Cov}
\DeclareMathOperator{\supp}{supp}\DeclareMathOperator{\Law}{Law}
\DeclareMathOperator{\Bern}{Bern}\DeclareMathOperator{\dist}{dist}
\newcommand{\E}{\mathbb E}\newcommand{\Pp}{\mathbb P}
\newcommand{\R}{\mathbb R}\newcommand{\Z}{\mathbb Z}
\newcommand{\ind}{\mathbf1}\newcommand{\ce}{C_{\mathrm E}}
\newcommand{\pe}{p_*}\newcommand{\qe}{q_*}\newcommand{\re}{r_*}
\newcommand{\te}{t_*}\newcommand{\be}{\beta_*}\newcommand{\nus}{\nu_*}
\newcommand{\phz}{\phi(0)}
\title{Small Lyapunov ratios and the structure\\of Berry--Esseen maximizers}
\author{Logan Bell}\date{6 October 2026}
\begin{document}\maketitle
\begin{abstract}
For changing finite independent centered arrays whose Lyapunov ratio
tends to zero and whose Berry--Esseen ratio tends to Esseen's constant,
we prove approximation by coordinatewise laws on one common lattice.
The sum of their squared and cubed Wasserstein errors is smaller than
the corresponding lattice scale, and the original additive zero-bias
transport bound is asymptotically an equality. If each row is an actual
global maximizer at its own finite horizon, then every original
representation is eventually exactly an IID two-point array padded by
zeros. The resulting Lyapunov threshold is qualitative. The proof
retains the absolute third moments of the original coordinates and
uses a relative conditional first-moment estimate to exclude phase
noise of arbitrarily small size.
\end{abstract}
\tableofcontents

\section{The results}
Let $X_1,\ldots,X_N$ be independent centered real variables with finite
third absolute moments. Write
\begin{equation}\label{eq:original-data}
 S=\sum_iX_i,\quad B^2=\sum_i\E X_i^2>0,\quad
 R=\sum_i\E|X_i|^3,\quad L=R/B^3,\quad
 C=L^{-1}d_K(S/B,Z),
\end{equation}
where $Z$ is standard normal. All sums of third moments below refer
to these original coordinates. Centered zero coordinates are allowed.
Set
\[
 \ce=\frac{3+\sqrt{10}}{6\sqrt{2\pi}},\quad
 \pe=\frac{4-\sqrt{10}}2,\quad\qe=1-\pe,\quad
 \te=\qe-\pe=\sqrt{10}-3,
\]
\[
 \re=\pe^2+\qe^2,\quad
 \be=\frac{\re}{\sqrt{\pe\qe}},\quad\nus=\re^{-1}.
\]
The variance-one Esseen law $W$ is specified by
$\Pp(W=\sqrt{\qe/\pe})=\pe$ and
$\Pp(W=-\sqrt{\pe/\qe})=\qe$.
For laws with finite $j$th moments, $W_j$ denotes Wasserstein distance,
with $W_j^j(X,Y)=\inf\E|\widetilde X-\widetilde Y|^j$, where the infimum
is over couplings with the stated marginal laws. The zero-biased law $S^*$ is defined by
$\E[Sf(S)]=B^2\E f'(S^*)$. The independent-sum zero-bias construction
and the sharp single-law bound give
\begin{equation}\label{eq:additive}
 \mathcal B(X):=\frac{2B^2W_1(S,S^*)}{R}\le1;
\end{equation}
see \citet{GoldsteinReinert1997,Goldstein2010,Tyurin2009} and the
equality classification in \citet[Theorem~1.1]{BellTransport2026}.

\begin{theorem}[Changing arrays near Esseen's constant]\label{thm:near}
In each row $n$, take a finite array satisfying
\eqref{eq:original-data}. Its length may change. Suppose
\begin{equation}\label{eq:near-assumption}
                         L_n\longrightarrow0,
                  \qquad C_n\longrightarrow\ce.
\end{equation}
There are positive gaps $h_n\sim R_n/(\re B_n^2)$ and centered laws
$X_{ni}^{\rm lat}$ on individual cosets of $h_n\Z$, coupled independently
coordinate by coordinate to the original laws, such that
\begin{equation}\label{eq:strong-lattice}
 \sum_iW_2^2(X_{ni},X_{ni}^{\rm lat})=o(h_n^2),\qquad
 \sum_iW_3^3(X_{ni},X_{ni}^{\rm lat})=o(h_n^3).
\end{equation}
Moreover $\mathcal B(X_n)\to1$.

There also exist subsets $I_n$ of active original coordinates,
$m_n=|I_n|$, and a common sign within each row, such that
$I_n$ retains $1-o(1)$ of both total variance and original price,
$m_nL_n^2/\be^2\to1$, and the padded IID laws
\[
 Q_{ni}=\Law(B_nW/\sqrt{m_n})\quad(i\in I_n),\qquad
 Q_{ni}=\delta_0\quad(i\notin I_n)
\]
obey
\[
 \sum_iW_2^2(X_{ni},Q_{ni})=o(B_n^2),\qquad
 \sum_iW_3^3(X_{ni},Q_{ni})=o(R_n).
\]
Here $W$ is the variance-one Esseen two-point law or its reflection.
The total variance of the target is exactly $B_n^2$ and its original
price is asymptotic to $R_n$.
\end{theorem}

The lattice comparators in \eqref{eq:strong-lattice} may have more
than two atoms. Their exact coordinate variances and cubic prices can
differ from those of the originals. The IID approximation has a weaker
error scale. A rare integer excursion can have a negligible contribution
relative to $R_n$ while preventing a binary approximation at scale $h_n$.

Let $C_N$ be the supremum of \eqref{eq:original-data} over all arrays
with at most $N$ coordinates. An actual global maximizer means an array
attaining this full supremum, with all centered finite-third-moment
laws admitted as competitors.

\begin{theorem}[Exact structure of small-Lyapunov globals]\label{thm:global}
Suppose each row is an actual global maximizer at its own finite horizon
$N_n$, and $L_n\to0$. Then every sufficiently late original array has
identically distributed nonzero two-point coordinates and zero
coordinates. The assertion applies to every original representation.
Equivalently, there exists an absolute $\varepsilon_0>0$ such that every
actual full global maximizer with $L<\varepsilon_0$ is padded IID binary.
\end{theorem}

The threshold in Theorem~\ref{thm:global} is ineffective. Zero padding
increases the horizon without changing $L$, so the theorem is stated
in terms of $L$. It does not specify the winning active count or binary
probability. The universal independent sharp-constant problem outside
the small-Lyapunov regime remains open.

The quantitative value theorem of \citet{BellSmall2026} concerns every
array with $L\le\exp(-\exp(97))$. Its penalized counterexample selection
and explicit local comparison are separate from the present classification
of actual globals, whose finite ratios can be below $\ce$. We use the
published asymptotic upper estimate
\begin{equation}\label{eq:asymptotic-upper}
 d_K(S/B,Z)\le\ce L+0.2538L^{4/3}\qquad(L\le10^{-3})
\end{equation}
from \citet[Corollary~4.18, p.~303]{Shevtsova2012}, rather than the
quantitative companion theorem, for the near-Esseen selection.
The moment and Fourier estimates extracted from that source are proved
with their retained losses in Section~\ref{sec:shape}.
The contamination contact argument follows
\citet[Section~2.1, equations~(7)--(12)]{HeChengCurrent2026}, applied
to one original coordinate rather than a common IID law.
The changing-array common-scale and lattice arguments, and the relative
conditional estimates required for exact phase collapse, are supplied
in the proofs below.

For exact structure, an approximation with an absolute error tending
to zero is insufficient. The remaining phase noise can be smaller than
that error by an arbitrary amount. We first confine every support phase,
then estimate conditional phase moments relative to the actual norm
$\E|T|$. Here a phase means the residual modulo the common lattice gap.
The integer count labels record the nearest lattice sites; the centered
rounding residuals may depend on their labels within each coordinate,
while the coordinate pairs remain independent. The integer endpoint
competitor has a loss of the actual phase norm, and
each comparison error is $o(\E|T|)$. This is the part of the argument
that uses full globality.

\section{Finite horizons and selected endpoints}\label{sec:finite}
We use closed and strict CDFs,
$F_S(y)=\Pp(S\le y)$ and $F_S^-(y)=\Pp(S<y)$. Continuity of the
normal target gives
\begin{equation}\label{eq:branches}
 d_K(S/B,Z)=\sup_y\max\{F_S(y)-\Phi(y/B),\,
                                  \Phi(y/B)-F_S^-(y)\}.
\end{equation}
Reflection turns the strict negative branch into a closed positive
branch. For positive discrepancy, the selected closed branch attains
its maximum: it is upper semicontinuous and tends to zero in both tails.

\begin{lemma}[Attainment at a fixed horizon]\label{lem:attainment}
For each finite $N$, $C_N$ is attained with zeros allowed. Every sequence
with ratio tending to $C_N$, normalized to total variance one, has a
subsequence converging coordinatewise in $W_3$ to a maximizing array.
\end{lemma}
\begin{proof}
A symmetric binary coordinate padded with zeros gives a positive
benchmark $c_0=\Phi(1)-1/2$. Since discrepancy is at most one, a ratio
at least $c_0$ implies $R\le1/c_0$. In the product weak topology the
centered $N$-tuples with total variance one and that cubic budget form
a compact set: tightness follows from the variance bound, and bounded
third moments give uniform integrability of first and second moments.
The cubic budget is lower semicontinuous. Convolution is weak-continuous
for fixed $N$. CDF bracketing against the continuous normal target shows
that discrepancy is weak-continuous. Finally
$R\ge\sum_i(\E X_i^2)^{3/2}\ge N^{-1/2}$, so discrepancy divided by
$R$ is upper semicontinuous. Arrays outside this compact set have smaller
ratio than the benchmark.

Along an extremizing subsequence, the discrepancies converge. Their
cubic sums therefore converge to the discrepancy of the limiting array
divided by $C_N$. Lower semicontinuity and maximality force equality
with its actual cubic sum. Each individual cubic moment converges,
because there are finitely many nonnegative terms. Weak convergence
together with third-moment convergence is $W_3$ convergence.
\end{proof}

The same supremum is obtained if every individual variance is required
to be positive, by replacing zero laws with arbitrarily small centered
bounded laws; that restricted class need not attain it. A separate
probability-Hessian result gives at most four support points in every
original global law and some equally optimal three-point representative
preserving each variance and price
\cite[Theorem~3.2]{BellExtremizers2026}.
The exact proof below works on the original laws directly and does not
use an existential replacement to change their support.

\section{Moment losses and a common scale}\label{sec:shape}
This section proves the approximate core used in both main theorems.
We normalize $B=1$ and omit row indices. Thus $R=L$. Put
$v_i=\E X_i^2$, $\sigma_i=\sqrt{v_i}$, $\rho_i=\E|X_i|^3$,
$\kappa_i=\E X_i^3$, and $b_i=\E|X_i|$. Define
\[
 \delta_i=\sqrt{10}\rho_i-|\kappa_i|-3b_iv_i.
\]

\begin{lemma}[Retained losses]\label{lem:retained}
Under \eqref{eq:near-assumption},
\begin{equation}\label{eq:retained-losses}
 \sum_i\delta_i=o(R),\quad
 \sum_i\E[|X_i|^3\ind_{\{|X_i|>R^{1/6}\}}]=o(R),\quad
 \sum_i|\kappa_i|-\Big|\sum_i\kappa_i\Big|=o(R),
\end{equation}
and $\sum_i\sigma_i^3/R\to1/\be$. After one reflection of a row,
let $P_i^*$ be the Esseen two-point law of variance $v_i$ with the
same positive orientation for every $i$. Then
\begin{equation}\label{eq:shape-W3}
                  \sum_iW_3^3(X_i,P_i^*)=o(R).
\end{equation}
\end{lemma}
\begin{proof}
First we record the equality and stability part of the scalar moment
inequality. Put $s=\sqrt{\pe\qe}$, $a=\sqrt{\qe/\pe}$,
$b=\sqrt{\pe/\qe}$ and $t=\te$. For a centered variance-one law
oriented to have nonnegative third moment, set
\[
 Q(x)=-6sx^2+6tx+6s,\qquad
 f(x)=x^3-\sqrt{10}|x|^3+3|x|.
\]
Direct factorization gives
\[
 Q(x)-f(x)=
 \begin{cases}
 (\sqrt{10}-1)(x-a)^2(x+2s),&x\ge0,\\
 (\sqrt{10}+1)(-x-b)^2(-x+2s),&x\le0.
 \end{cases}
\]
Since $\E Q(X)=0$, $\delta(X)=\E(Q-f)\ge0$; equality forces
the two Esseen support points, with their centered probabilities.
This also checks the equality statement needed from
\citet[Theorem~3.1, pp.~256--262]{Shevtsova2012}.

For $c_0=2\pe(\sqrt{10}-1)>0$, each factor is at least $c_0$
times the cube of the distance to its half-line contact. Round a
variance-one law to $a$ on the positive half-line and to $-b$ on the
negative half-line. If $e$ is its cubic rounding cost, then
$e\le\delta(X)/c_0$ and centering gives
$|\Pp\{\hbox{rounded value}=a\}-\pe|\le e^{1/3}/(a+b)$.
Couple the rounded law to the exact Esseen probabilities. The inequality
$(u+v)^3\le4(u^3+v^3)$ yields
\[
 W_3^3(X,W_+)\le4e+4(a+b)^2e^{1/3}.
\]
After rescaling coordinates and applying H\"older,
\begin{equation}\label{eq:shape-stability}
 \frac{\sum_iW_3^3(X_i,P_i^{\rm oriented})}{R}
 \le\frac4{c_0}\bar\delta+
       4(a+b)^2c_0^{-1/3}\bar\delta^{1/3},\qquad
 \bar\delta=R^{-1}\sum_i\delta_i.
\end{equation}
Here the orientations initially follow the signs of individual
$\kappa_i$; zero laws have zero cost.

The proof on pp.~289--296 of \citet{Shevtsova2012} permits retention
of its nonnegative losses. We give the details of this extraction.
Put $U=R^{1/6}$, $Y_i=X_i\ind_{\{|X_i|\le U\}}$,
$Z_i=Y_i-\E Y_i$, and $T_i=\E[|X_i|^3\ind_{\{|X_i|>U\}}]$.
The source's $I_{14},I_{15},I_{16}$ comparisons, with its truncation
parameter $u=1$, and Lemma~4.3 on pp.~272--274 give
\[
 \Delta\le\phz(I_{16}+R/2)+O(R^{7/6}),
\]
\[
 I_{16}=\frac{\sqrt{10}R}{6}-\frac16\sum_i\delta(Z_i)
                   +I_{17}-I_{18},
\]
\[
 I_{17}=\frac{\sqrt{10}}6\sum_i(\E|Z_i|^3-\E|Y_i|^3),\qquad
 I_{18}=\frac{\sqrt{10}-1}{6}\sum_iT_i.
\]
The second identity is the moment decomposition before discarding
its deficit. At the parameter $\sqrt{10}$ the source's moment
remainder is zero. All remaining terms are of order $R^{7/6}$ or
smaller, uniformly over the independent row.

Let $b_{0i}=\rho_i^{1/3}/U\le R^{1/6}$. The bound
$|\E Y_i|\le T_i/U^2$ and cubic Taylor expansion show
\[
 \E|Z_i|^3-\E|Y_i|^3
       \le(3b_{0i}^2+3b_{0i}^4+b_{0i}^6)T_i.
\]
For $R\le1/4096$, $b_{0i}\le1/4$ and the bracket is less than
$1/3$. Since $\sqrt{10}/3<(\sqrt{10}-1)/2$, the displayed
decomposition implies
\[
 \Delta\le\ce R-\frac{\phz}{6}\sum_i\delta(Z_i)
            -\frac{\phz(\sqrt{10}-1)}{12}\sum_iT_i+O(R^{7/6}).
\]
Near-Esseen competition forces both losses to be $o(R)$. Cubic Taylor
bounds give
$|\rho_i-\rho(Z_i)|+|\kappa_i-\kappa(Z_i)|\le K T_i$, while
$\E|Z_i|\Var Z_i-\E|X_i|\Var X_i\le b_{0i}^2T_i$.
Consequently $\delta_i\le\delta(Z_i)+K'T_i$, proving the first two
claims in \eqref{eq:retained-losses} and the individually oriented
version of \eqref{eq:shape-W3}.

To retain orientation we use the aggregate low-frequency integral,
before the source's individual cubic triangle inequality. Write
$K_3=\sum_i\kappa_i$, $\nu=1+\sum_i b_iv_i/R\in[1,2]$ and
$T=2\pi/(\nu R)$. On each fixed compact frequency interval the
just-proved tail estimate gives
\[
 \chi(u)=e^{-u^2/2}(1-iK_3u^3/6)+o(R),
\]
uniformly there. Indeed the summed Taylor remainders are at most
$K[UR|u|^4+\sum_iT_i|u|^3]=o(R)$ and the summed logarithmic
quadratic errors are $O(\sum_i v_i^2+\sum_i\rho_i^2)=o(R)$.

We use Lemma~4.1 and (4.4)--(4.7) on pp.~263--267 of the source.
Its smoothing kernel is
\[
 K(t)=\frac{1-|t|}{2}+\frac{i}{2}
       \left\{(1-|t|)\cot(\pi t)+\frac{\operatorname{sign}t}{\pi}\right\},
                                      \qquad0<|t|<1.
\]
The constants $\kappa_1$ and $t_1\in(1/2,1)$ below are the universal
constants in the source's (4.4)--(4.6); in particular
$t^2(1-4\pi\kappa_1t)$ is strictly positive on $0<t\le t_1$.
Choose a fixed small $t_0>0$ so that its deleted-product bound (4.7)
has Gaussian exponent at least $1/4$ for $0\le u\le t_0T$.
The one-coordinate second-order Taylor bound against a same-variance
normal gives
\[
 R^{-1}|\chi(u)-e^{-u^2/2}|\le K|u|^3e^{-u^2/4}
                           \quad(0\le u\le t_0T).
\]
The Prawitz kernel satisfies
$2|K(u/T)|/T\le1.0253/(\pi u)$ and converges to $1/(\pi u)$.
Dominated convergence therefore evaluates the source's low integral as
$\phz|K_3|/6+o(R)$. Its middle interval is exponentially small:
the function $t^2(1-4\pi\kappa_1t)$ is positive on the chosen
$[t_0,t_1]$. At the final endpoint $|K(1-u)|\le Ku$ and
$1-\cos(2\pi u)\ge8u^2$, giving $O(T^{-2})$. The Gaussian endpoint
terms are at most $\sqrt{\pi/2}/T+O(T^{-2})$. Thus
\begin{equation}\label{eq:signed-loss}
 \Delta\le\phz\left\{\frac R2+\frac12\sum_i b_iv_i
                                  +\frac{|K_3|}{6}\right\}+o(R)
 =\ce R-\frac{\phz}{6}\left\{\sum_i\delta_i+
                  \sum_i|\kappa_i|-|K_3|\right\}+o(R).
\end{equation}
This proves the orientation loss. Reflect so $K_3\ge0$. Moment
continuity under the cubic couplings gives
\[
 \sum_i|\kappa_i-\kappa(P_i^{\rm oriented})|
 \le3\left(R+\be\sum_i\sigma_i^3\right)^{2/3}
                 \left(\sum_iW_3^3(X_i,P_i^{\rm oriented})\right)^{1/3}
 =o(R).
\]
The negatively oriented target coordinates then have
$\sum\sigma_i^3=o(R)$; reflecting just those targets costs $o(R)$
in cubic transport. This proves the common-orientation assertion.

Finally \citet[Theorem~4.14, pp.~300--303]{Shevtsova2012} gives, for
each fixed $c>2\phz/3$,
\[
 \Delta\le cR+K(c)\sum_i\sigma_i^3+O_c(R^{4/3}),
\]
where its Bernoulli envelope is
\[
 K(c)=\sup_{0<p\le1/2}
 \frac{\phz(3+1-2p)/6-c\{p^2+(1-p)^2\}}{\sqrt{p(1-p)}}.
\]
At $c=\ce$ its unique maximizing parameter is $\pe$,
$K(\ce)=0$, and $K'(\ce)=-\be$. The numerator's exact maximum
and uniqueness follow either by differentiation or from the sharp
moment equality above; parameters near zero are separated since
$c>2\phz/3$. Any subsequential limit $z\in[0,1]$ of
$\sum\sigma_i^3/R$ satisfies $\ce\le c+K(c)z$ on both sides of
$\ce$. Differentiation there forces $1-\be z=0$.
\end{proof}

\begin{lemma}[An actual Fourier window]\label{lem:window}
For every fixed $\epsilon>0$, a sufficiently small fixed $t_0>0$ gives
\[
 T'=(1+\epsilon)\frac{2\pi}{\nu R},\qquad
 H=2\int_{t_0}^1|K(u)|\,|\chi(T'u)|\,du,
\]
with $\nu\to\nus$ and
\begin{equation}\label{eq:window}
 \liminf H/R\ge\frac{\phz\nus\epsilon}{2(1+\epsilon)}>0.
\end{equation}
For one fixed such window there are $\eta>0$ and $e_n\to0$ such that
simultaneously for every subset of coordinates $A$,
\begin{equation}\label{eq:subset-inverse}
             \sum_{i\in A}v_i\le\eta^{-1}
                          \frac{\sum_{i\in A}\rho_i}{R}+e_n.
\end{equation}
Moreover $\sum_iW_2^2(X_i,P_i^*)\to0$.
\end{lemma}
\begin{proof}
Common-orientation moment continuity from Lemma~\ref{lem:retained}
gives $\nu\to1+3/\sqrt{10}=\nus$ and $K_3/R\to1/\sqrt{10}$.
Repeat the source's Prawitz inequality at the enlarged cutoff $T'$,
choosing $t_0$ smaller so its low-frequency domination remains valid.
Keep the actual high integral $H$ rather than extending a previous
high-frequency upper estimate beyond its domain. The low integral is
$\phz|K_3|/6+o(R)$ and the Gaussian endpoint terms are at most
$\phz\nu R/[2(1+\epsilon)]+O(R^2)$. Near-Esseen competition gives
\eqref{eq:window}.

For any centered law and an independent copy,
\[
 |\chi_i(u)|^2\le1-v_iu^2+\tfrac43\rho_i|u|^3.
\]
This follows from $\cos x\le1-x^2/2+|x|^3/6$ and
$\E|X_i-X_i'|^3\le8\rho_i$. Independence gives for a subset
\[
 |\chi_A(u)|\le\exp(-V_Au^2/2+2R_A|u|^3/3).
\]
The fixed window lies in $[a/R,C/R]$ with $a>0$. Choose
$0<\eta\le3/(8C)$. If $R_A/V_A\le\eta R$, then the full product
on this window is at most $\exp(-a^2V_A/(4R^2))$.
The common lower bound $H\ge hR$ implies
$V_A\le e_n=(4R^2/a^2)\log(2\sup|K|/(hR))$.
Otherwise $V_A<\eta^{-1}R_A/R$. These two cases prove
\eqref{eq:subset-inverse} with one eventual index for all subsets.

Let $E_i=W_3^3(X_i,P_i^*)$, so $\sum E_i=o(R)$. Fix $0<\zeta<1$.
On the set $E_i>\zeta\rho_i$, the original price and, by
\eqref{eq:subset-inverse}, variance tend to zero. On its complement,
Minkowski gives
$\rho_i/\sigma_i^3\le\be/(1-\zeta^{1/3})^3$ and
\[
 W_2^2(X_i,P_i^*)\le E_i^{2/3}
 \le v_i\left\{\frac{\zeta\be}{(1-\zeta^{1/3})^3}\right\}^{2/3}.
\]
On the exceptional set a centered independent same-variance coupling
costs $2v_i$. Let the row tend to infinity, then let $\zeta\downarrow0$.
\end{proof}

\begin{proposition}[A variance-and-price-full Esseen core]\label{prop:core}
The IID target conclusion of Theorem~\ref{thm:near} holds. In gap units
$h=\nus R$, there is a further retained set of cardinality $M$ with
\begin{equation}\label{eq:raw-core}
 M R^2\to\be^2,\quad
 \sup_{i\in I}W_3(\Law(X_i/h),\Law(B_{\pe}-\pe))\to0,
\end{equation}
retaining $1-o(1)$ of original variance and price. In these units the
total original variance and price satisfy
$V\sim\pe\qe M$ and $P/V\to\re$, and its signed cubic sum divided
by $V$ tends to $\te$.
\end{proposition}
\begin{proof}
Consider the variance probability measures
$\mu_n=\sum_i v_i\delta_{\sigma_i/R}$.
Their first moments tend to $s_0=1/\be$, so they are tight at
infinity. They have no limiting mass at zero. To prove this, discard
the set $E_i>\zeta\rho_i$ as in the last lemma. On the remaining set
$\rho_i/\sigma_i^3\le K_\zeta$. For $\sigma_i\le\delta R$, its
price is at most $K_\zeta\delta R$ times its variance. The subset
inverse gives limiting variance at most $\eta^{-1}K_\zeta\delta$.
Portmanteau on the relative-open interval $[0,\delta)$ followed by
$\delta\downarrow0$ excludes a zero atom from every limit $\mu$.
Fatou gives $\int z\,d\mu\le s_0$; first-moment uniform
integrability has not yet been assumed.

At each fixed $\epsilon>0$, \eqref{eq:window} selects a frequency
$u_n\in[t_0T',T']$ with $|\chi(u_n)|\ge c_\epsilon R$.
Let $w_n=Ru_n$. It lies in a compact positive interval with
$\limsup w_n\le2\pi(1+\epsilon)/\nus$. Independence gives
$\sum_i(1-|\chi_i(u_n)|^2)=O_\epsilon(\log(1/R))$.
The same-variance binary target has defect
$4\pe\qe\sin^2((a+b)\sigma_i u_n/2)$.
On $\sigma_i\le KR$, the original variance-weighted defects tend
to zero, and their transfer to target defects is at most
\[
 2|u_n|\Big(\sum_{\sigma_i\le KR}v_i^2\Big)^{1/2}
                          \Big(\sum_iW_2^2(X_i,P_i^*)\Big)^{1/2}=o(1).
\]
The untruncated target contribution is bounded by $\mu_n([K,\infty))$.
First let the row tend to infinity, then $K\to\infty$. Along a weak
subsequence with $w_n\to w_\epsilon$, this proves
\[
          \int\sin^2((a+b)zw_\epsilon/2)\,d\mu(z)=0.
\]
The limit is supported on positive integer multiples of
$2\pi/((a+b)w_\epsilon)$, hence on $[s_0/(1+\epsilon),\infty)$;
we used $\be/(a+b)=1/\nus$. Repeat for $\epsilon=1/j$, taking
further subsequences of the same weak limit. That limit is supported
on $[s_0,\infty)$ and has mean at most $s_0$, so it is $\delta_{s_0}$.
The exact first moments converge; therefore
\[
 \int|z-s_0|\,d\mu_n
   =\int z\,d\mu_n-s_0+2\int(s_0-z)_+\,d\mu_n\longrightarrow0.
\]

Trim a slowly shrinking band about $s_0$ and denote its count by $m$.
It retains total variance $1-o(1)$ and target cubic resource
$\be\sum_{I}\sigma_i^3=R(1-o(1))$. The original price has the
same retention, since uniformly over subsets
\[
 \sum_i|\rho_i-\be\sigma_i^3|
 \le3\left(R+\be\sum_i\sigma_i^3\right)^{2/3}
                       \left(\sum_iE_i\right)^{1/3}=o(R).
\]
Counting retained variances gives $mR^2/\be^2\to1$.
Replace the core's individually scaled targets by $W_+/\sqrt m$.
Their summed squared scale error tends to zero, and their summed
cubed scale error is $o(R)$; both follow by expanding
$\sigma_i-1/\sqrt m=R(\sigma_i/R-s_0)+R(s_0-1/(R\sqrt m))$.
The outside-to-zero costs are exactly its variance and original price.
This proves the stated IID conclusion after restoring $B$.

In gap units $h=\nus R$, the core's total cubic transport error is
$o(M)$. Remove indices whose error in gap units exceeds a positive number tending
slowly to zero. Their count is $o(M)$ and the preceding uniform moment
continuity estimate preserves variance and price retention. This gives
\eqref{eq:raw-core}; its moment limits follow by uniform $W_3$ continuity.
\end{proof}

\section{A first nonzero Fourier peak and integer lifts}\label{sec:alias}
\begin{lemma}\label{lem:first-alias}
At a gap asymptotic to that in Proposition~\ref{prop:core}, the original
variables in gap units have representations
\[
 Y_i=J_i+Z_i,\qquad J_i=K_i-\E K_i,\quad K_i\in\Z,\quad
 \E Z_i=0,\quad |Z_i|\le1.
\]
The pairs are independent between coordinates. On the retained core
$K_i$ converges uniformly in $W_3$, after integer translations, to
$\Bern(\pe)$. Outside it the summed centered integer variance and
cubic price are $o(M)$. Writing $K=\sum K_i$ and $T=\sum Z_i$,
\begin{equation}\label{eq:initial-lifts}
 Q=\Var T=O(1),\qquad \Cov(K,T)=o(\sqrt M),\qquad
 A=\Var K\sim\pe\qe M.
\end{equation}
\end{lemma}
\begin{proof}
Fix $0<\epsilon<1/2$ in Lemma~\ref{lem:window}. The frequency in gap units
window stays in a compact subinterval of $(0,4\pi)$ containing only
the alias $2\pi$, and its actual characteristic-modulus area is
at least $c/\sqrt M$.
Uniform gap-unit $W_3$ convergence in \eqref{eq:raw-core} gives uniform
convergence of every core characteristic function and its first two
derivatives on that window. Polynomial-factor differences are controlled
by the $L^2$ coupling and the bounded Bernoulli target, so no fourth
moment is required. Around $2\pi$ those factors are nonzero and their
log-modulus second derivatives are at most $-c$. Away from a fixed
neighborhood the retained product is exponentially small.

Let $s_0\to2\pi$ be the exact retained-product modulus maximum.
Its product is bounded by $\exp(-cM(s-s_0)^2)$ near that point.
The area lower bound therefore yields a full-product point
$s_1=s_0+O(M^{-1/2})$ of fixed positive height. The centered omitted
original sum has variance $o(M)$, so its characteristic function changes
by $o(1)$ from $s_1$ to $s_0$:
\[
 |\chi_{\rm rem}(s_1)-\chi_{\rm rem}(s_0)|
 \le |s_1-s_0|\sqrt{\Var Y_{\rm rem}}=o(1).
\]
The full product at the exact core maximum also has fixed positive height.
Rescale the original variables by $s_0/(2\pi)$.

Align each marginal characteristic argument at $2\pi$ and write
$Y_i=\theta_i+K_i+E_i$ with $E_i\in[-1/2,1/2]$, using a fixed
nearest-integer tie rule. Then
\[
 \E\sin(2\pi E_i)=0,\quad \E\cos(2\pi E_i)=r_i>0,\quad
 \sum_i\E E_i^2\le\tfrac18\sum_i(1-r_i)=O(1).
\]
Centering gives $Z_i=E_i-\E E_i$ and the claimed representation.
The exact core logarithmic derivative is zero, hence
$\sum_{I}\E[Y_i\sin(2\pi E_i)]/r_i=0$.
On the core $r_i\to1$ uniformly. Removing these denominators costs
$o(1)$ by Cauchy and the bounded sum of defects. Replacing the sine
by $2\pi E_i$ costs $O(1)+o(\sqrt M)$: couple $Y_i$ to its bounded
binary target and use
\[
 \sum_I\E[|Y_i||E_i|^3]
 \le K\sum_I\E|E_i|^3+
       \left(\sum_I\E|Y_i-W_i|^2\right)^{1/2}
       \left(\sum_I\E|E_i|^6\right)^{1/2}.
\]
The omitted covariance is $o(\sqrt M)$ by its original variance and
bounded phase variance. Thus $\Cov(\sum Y_i,T)=o(\sqrt M)$ and
$\Cov(K,T)=\Cov(\sum Y_i,T)-Q=o(\sqrt M)$.

Core rounding gives uniform $W_3$ convergence of $K_i$ to $\Bern(\pe)$:
outside a fixed small coupling-error event its two nearest integers are
exactly the binary labels, while the remaining cubic rounding error
is bounded by a constant times the cubic target-coupling error.
For the omitted laws, $J_i=Y_i-Z_i$ and $|Z_i|\le1$ imply
\[
 \sum_{\rm rem}\Var K_i\le2\sum_{\rm rem}\Var Y_i+2Q=o(M),\quad
 \sum_{\rm rem}\E|J_i|^3\le4\sum_{\rm rem}\E|Y_i|^3+4Q=o(M).
\]
All prices in the last expression remain separate original-coordinate
prices. These estimates also give $A\sim\pe\qe M$.
\end{proof}

\section{Joint integer and residual limit theorems}\label{sec:joint}

Under the decomposition $Y_i=J_i+Z_i$ of Lemma~\ref{lem:first-alias}, write
$J_i=K_i-\E K_i$, $K=\sum_{i=1}^N K_i$, $\mu=\E K$, $T=\sum_{i=1}^N Z_i$,
$A=\Var K$, $Q=\Var T=O(1)$, $S=K-\mu+T$, and $V=\Var S$. Because the core
$I$ of size $M$ converges uniformly in $W_3$ to $\Bern(p)$ ($p=\pe$) and
carries all but $o(M)$ of the integer variance and third absolute moments,
while $\Cov(K,T)=o(\sqrt M)$, we have
\begin{equation}\label{eq:joint-moments}
 A\sim p(1-p)M,\qquad V=A+o(\sqrt A),\qquad
 \gamma=\frac{\sum_{i=1}^N\E J_i^3}{A}\longrightarrow1-2p.
\end{equation}

\begin{lemma}[Global Edgeworth and local estimates for the integer sum]\label{lem:integer}
Uniformly over all $j\in\Z$, with $x_j=(j-\mu)/\sqrt A$,
\begin{align}
 \Pp(K\le j)&=\Phi(x_j)+\frac{\phi(x_j)}{\sqrt A}
             \left\{\frac12+\frac\gamma6(1-x_j^2)\right\}
                        +o(A^{-1/2}),\label{eq:integer-edgeworth}\\
 \sqrt A\,\Pp(K=j)&=\phi(x_j)+o(1),\label{eq:integer-local}
\end{align}
and uniformly over all $k\in\Z$,
\begin{equation}\label{eq:all-layer}
                  \E\bigl[T^2\ind_{\{K=k\}}\bigr]\le \frac{K_0 Q}{\sqrt A}.
\end{equation}
All three estimates remain uniform after deleting any fixed number of
summands.
\end{lemma}
\begin{proof}
On the core $I$, every $K_i$ has atom probabilities at $0$ and $1$ bounded
below, so after deleting a bounded number of summands the remaining product
satisfies the torus damping bound
\begin{equation}\label{eq:torus-damping}
        |\chi_{\rm deleted}(u)|\le K\exp\bigl(-cM\sin^2(u/2)\bigr)
                                                   \qquad(|u|\le\pi).
\end{equation}
Expanding $\log\chi_I(u)=-Wu^2/2-\ii\kappa_Iu^3/6+O(M(\epsilon_n|u|^3+|u|^4))$
near $u=0$, inverting against the Gaussian--Hermite lattice sequence via
Euler--Maclaurin summation, and convolving with the outside block (which has
variance and third-moment sum $o(A)$) yields \eqref{eq:integer-edgeworth}
and \eqref{eq:integer-local}. For \eqref{eq:all-layer}, expanding the Fourier
transform of $T^2\ind_{\{K=k\}}$ into diagonal terms (bounded in sum by $Q$)
and off-diagonal terms (with factors $|\E[Z_ie^{\ii uJ_i}]|\le|u|\sqrt{\Var(K_i)\E Z_i^2}$)
and integrating $KQ(1+Au^2)e^{-cAu^2}$ over $[-\pi,\pi]$ gives the uniform
bound $K_0Q/\sqrt A$.
\end{proof}

\begin{lemma}[Asymptotic independence and joint local limit]\label{lem:joint-local}
Along any subsequence for which $T\Rightarrow U$, we have
$((K-\mu)/\sqrt A,T)\Rightarrow(G,U)$, where $G\sim N(0,1)$ is independent
of $U$ and $\E U=0$. Moreover, for every integer sequence $k$ with
$(k-\mu)/\sqrt A\to x\in\R$,
\begin{equation}\label{eq:joint-local}
        \sqrt A\,\Pp(K=k,\ T\in du)\Longrightarrow\phi(x)\Law(U)(du).
\end{equation}
\end{lemma}
\begin{proof}
For fixed $u,v\in\R$ and $s=u/\sqrt A$, expanding
$\sum_{i\in I}\log(\E e^{\ii sJ_i+\ii vZ_i}/\E e^{\ii vZ_i})$ through second
order in $s$ and using $\Cov(K,T)=o(\sqrt A)$ gives
$\E e^{\ii s(K-\mu)+\ii vT}=e^{-u^2/2+o(1)}\E e^{\ii vT}$, proving joint
weak convergence. To prove \eqref{eq:joint-local}, invert the Fourier
transform in the first coordinate at $k$. Since $\sup_{i\in I}\E Z_i^2\to0$,
on a fixed small cell $|s|\le c_0$ the twisted core factors
$\E e^{\ii sJ_i+\ii vZ_i}$ converge uniformly in $C^2$ to the centered
$\Bern(p)$ factor, are bounded away from zero, and have log-modulus second
derivatives at most $-c$ while their log-modulus first derivatives at $s=0$
sum to at most $K_v\sqrt A$ by Cauchy--Schwarz. Thus
\[
 \left|\prod_{i\in I}\E e^{\ii sJ_i+\ii vZ_i}\right|
          \le\exp\bigl(K_v|s|\sqrt A-cAs^2\bigr)\qquad(|s|\le c_0),
\]
whereas on $c_0\le|s|\le\pi$ the uniformly binary-like twisted core product is
exponentially small in $M$. Substituting $s=w/\sqrt A$, dominated convergence
and tightness from \eqref{eq:all-layer} together with \eqref{eq:integer-local}
yield \eqref{eq:joint-local}.
\end{proof}

\begin{lemma}[Asymptotic threshold envelopes]\label{lem:endpoints}
If $y/\sqrt A\to x\in\R$, let $j=\lfloor y+\mu\rfloor$ and suppose along a
subsequence that $\theta=y+\mu-j\to\theta_0\in[0,1]$ and $T\Rightarrow U$.
Then
\begin{align}
 \limsup_{n\to\infty}\sqrt A\,\bigl\{F_S(y)-\Phi(y/\sqrt V)\bigr\}
 &\le\phi(x)\left\{\frac12+\frac{1-2p}{6}(1-x^2)
                         -\E\{\theta_0-U\}\right\},\label{eq:closed-envelope}\\
 \liminf_{n\to\infty}\sqrt A\,\bigl\{F_S^-(y)-\Phi(y/\sqrt V)\bigr\}
 &\ge\phi(x)\left\{-\frac12+\frac{1-2p}{6}(1-x^2)\right\},
                                                        \label{eq:strict-envelope}
\end{align}
where $\{w\}=w-\lfloor w\rfloor$ denotes fractional part. If
$|y|/\sqrt A\to\infty$, then $\sqrt A\,|F_S^\pm(y)-\Phi(y/\sqrt V)|\to0$.
\end{lemma}
\begin{proof}
Decomposing $F_S(y)-F_K(j)=\sum_{\ell\in\Z}\{\Pp(K=j+\ell,T\le\theta-\ell)-\ind_{\{\ell\le0\}}\Pp(K=j+\ell)\}$,
truncating at $|\ell|\le H$ via \eqref{eq:all-layer} (which bounds the tail
by $KQ/(H\sqrt A)$), applying \eqref{eq:joint-local} and Portmanteau on closed
half-lines, and sending $H\to\infty$ gives
$\phi(x)\E\lfloor\theta_0-U\rfloor=\phi(x)(\theta_0-\E\{\theta_0-U\})$ (since
$\E U=0$). Combining this with \eqref{eq:integer-edgeworth} at $j=y+\mu-\theta$
proves \eqref{eq:closed-envelope}, and open half-lines similarly prove
\eqref{eq:strict-envelope}.
\end{proof}

\subsection{Modular concentration and proof of Theorem~\ref{thm:near}}

For a sequence satisfying \eqref{eq:near-assumption} with $p=\pe$, we have
$P/V\to\re$ and $\kappa/V\to\te$ by Proposition~\ref{prop:core}. The negative
strict branch \eqref{eq:strict-envelope} has normalized upper limit
$\phz(1/2-\te/6)/\re<\ce$, while the positive closed branch
\eqref{eq:closed-envelope} has unique maximum $\ce$ at $x=0$ and
$\E\{\theta_0-U\}=0$. Thus $U$ is concentrated on the integer coset
$\theta_0+\Z$, so there exist phases $\vartheta_n$ such that
\begin{equation}\label{eq:modular}
              \E\dist(S-\vartheta_n,\Z)\longrightarrow0.
\end{equation}

To deduce $\mathcal B(X_n)\to1$ from \eqref{eq:modular}, let $q=\qe$,
$t=2q-1=\te$, and let $P_q(u)$ be the $1$-periodic sawtooth function equal to
$-2(1-q)u$ on $[0,q]$ and $2q(u-1)$ on $[q,1]$. Then $P_q(x-\vartheta)$ has
mean $-q(1-q)$ on the circle and vanishes on $\vartheta+\Z$ (with
$|P_q(x-\vartheta)|\le2\dist(x-\vartheta,\Z)$), while the combined test
function $-tx+P_q(x-\vartheta)$ is $1$-Lipschitz on $\R$. Using
$\E S^*=\kappa/(2V)$ and $|\E e^{\ii uS^*}|\le1/(|u|\sqrt V)$ in the Fourier
series of $P_q(x-\vartheta)$ gives, by Kantorovich--Rubinstein duality,
\begin{equation}\label{eq:modular-dual}
 W_1(S,S^*)\ge\frac{t\kappa}{2V}+q(1-q)
  -2\E\dist(S-\vartheta,\Z)-\frac{\zeta(3)}{\pi^3\sqrt V}.
\end{equation}
Taking $\liminf$ via \eqref{eq:modular} gives
$\liminf W_1(S,S^*)\ge\te^2/2+\pe\qe=\re/2$, which matches the upper bound
$W_1(S,S^*)\le P/(2V)\to\re/2$ and proves $\mathcal B(X_n)\to1$.

Finally, \eqref{eq:modular} implies $|\E e^{2\pi\ii S}|=\prod_i|\E e^{2\pi\ii Z_i}|\to1$,
so $Q=\sum_i\E Z_i^2\le\frac14\sum_i(1-|\E e^{2\pi\ii Z_i}|^2)\to0$. Since
$|Z_i|\le1$, the direct couplings $(Y_i,J_i)$ satisfy
$\sum_i W_2^2(Y_i,J_i)\le Q\to0$ and $\sum_i W_3^3(Y_i,J_i)\le Q\to0$.
Restoring the gap scale $h_n$ proves \eqref{eq:strong-lattice} and completes
the proof of Theorem~\ref{thm:near}.

\begin{remark}
While $\mathcal B(X_n)\to1$ whenever $L_n\to0$ and $C_n\to\ce$, the pointwise
inequality $C\le\ce\mathcal B(X)$ does not hold for general distributions.
Indeed, for $N=1$, replacing the standard normal law $Z$ on $[-a,a]$ by three
atoms at $\{-a,0,a\}$ with matching mean and second moment (weights
$Pr/2,P(1-r),Pr/2$ where $P=\Pp(|Z|\le a)$ and $r=\E[Z^2\ind_{\{|Z|\le a\}}]/(a^2P)\to1/3$)
gives $d_{\mathrm K}=P(1-r)/2$ and $W_1=(aP/2)\{(1-r)^2+r^2\}$, so
$C/\mathcal B(X)=(1-r)/[2a\{(1-r)^2+r^2\}]\sim3/(5a)\to\infty$ as $a\downarrow0$.
\end{remark}

\section{Variational contacts and uniform phase confinement}\label{sec:contacts}

For the remainder of the paper, each array in the sequence is a global
maximizer of $C_{N_n}$ at its number of summands $N_n$ and satisfies
$L_n\to0$. Since the number of nonzero summands is at least $L_n^{-2}\to\infty$,
we have $N_n\to\infty$, and comparing with an $N_n$-summand i.i.d.\ Esseen
array gives $C_{N_n}\to\ce$, so Theorem~\ref{thm:near} applies.

\begin{lemma}[Variational contact and support coercivity]\label{lem:coercivity}
Normalize total variance to $B=1$, so $R=L$. For all sufficiently late global
maximizers, every support point of every summand satisfies
$|u|\le246R^{1/3}$. At an attaining positive closed threshold $z$, with
$D=F(z)-\Phi(z)$, $\alpha=\phi(z)$, $A_0=(3D+z\alpha)/2$, $v=\E X_i^2$,
$\rho=\E|X_i|^3$, and $M_i=\E[X_i|X_i|]$, every summand $X_i$ satisfies
\begin{equation}\label{eq:full-contact-unit}
 F_{-i}(z-u)+(\alpha+3CM_i)u+A_0u^2-C|u|^3
       \le F(z)+A_0v-C\rho\qquad(u\in\R),
\end{equation}
with equality at every $u\in\supp X_i$.
\end{lemma}
\begin{proof}
Perturbing $\Law(X_i)$ to $(1-e)\Law(X_i)+e\delta_u$ at the fixed closed
threshold $z$ and recentering and restandardizing gives \eqref{eq:full-contact-unit}
by maximality (adapting \citet[Section~2.1, eqs.~(7)--(12)]{HeChengCurrent2026}
to a single coordinate), and upper semicontinuity of $F_{-i}(z-u)$ extends
equality from a set of full $\Law(X_i)$-measure to $\supp X_i$. Substituting
the deleted Berry--Esseen bound $|F_{-i}(z-u)-\Phi((z-u)/\sqrt{1-v})|\le(123/10)R$
\citep[Theorem~2.1]{ChenShao2001} into the equality at $u=R^{1/3}w\in\supp X_i$
gives $(99/500)w^3\le1023/40+(129/10)w+(203/20)w^2$, which forces $w<246$.
\end{proof}

In the gap units of Lemma~\ref{lem:first-alias}, write $Y_i=J_i+Z_i$ with
$V=\sum_i\E Y_i^2$, $d=V^{-1/2}$, $P=\sum_i\E|Y_i|^3$, $Q=\sum_i\E Z_i^2$,
$T=\sum_i Z_i$, and standardized threshold $z=y/\sqrt V$. By
Theorem~\ref{thm:near} and Lemma~\ref{lem:coercivity}, $|Y_i|\le H=O(V^{1/3})$,
$\sup_i\E|Y_i|^3=o(V)$, and
\begin{equation}\label{eq:contact-inputs}
 \frac PV\to\re,\qquad \frac\kappa V\to\te,\qquad Q\to0,\qquad
 z\to0,\qquad \frac Dd\to\phz\left(\frac12+\frac\te6\right),\qquad
 \Cov(K,T)=o(\sqrt V).
\end{equation}

\begin{lemma}[Uniform support-phase confinement]\label{lem:whole-phase}
For the decomposition $Y_i=J_i+Z_i$,
\begin{equation}\label{eq:whole-phase}
       \eta_n:=\max_{1\le i\le N_n}\|Z_i\|_\infty\longrightarrow0,
\end{equation}
and the fractional part $\{y+\mu\}$ converges to $0$ modulo $1$. In
particular, whenever $t:=\E|T|>0$, we have $Q=o(t)$.
\end{lemma}
\begin{proof}
For any deleted coordinate $i$, put $v_i=\E Y_i^2$, $\rho_i=\E|Y_i|^3$,
$w_i=V-v_i$, and $\gamma_i=\kappa(Y_{-i})/w_i$, where $\kappa(Y_{-i})$ is the
deleted third-moment sum. For $|u|\le H$, if the deleted fractional part
$\tau_i(u)=\{y-u+\mu_{-i}\}$ lies in $[\epsilon,1-\epsilon]$, then
Lemma~\ref{lem:integer} applied to the deleted sum gives
\begin{equation}\label{eq:away-phase}
               F_{-i}(y-u)=G_i(u)-d\phz\tau_i(u)+o(d),
\end{equation}
where $G_i(u)=\Phi(x_i)+w_i^{-1/2}\phi(x_i)\{1/2+\gamma_i(1-x_i^2)/6\}$
($x_i=(y-u)/\sqrt{w_i}$) is the smooth deleted Edgeworth approximation. In
gap units, the contact score and level from \eqref{eq:full-contact-unit} are
\[
 g_i(u)=F_{-i}(y-u)+b_iu+A_0u^2-\lambda|u|^3,\qquad
 \ell_i=F(y)+A_0v_i-\lambda\rho_i,
\]
where
\[
 b_i=d\phi(z)+3\lambda\E[Y_i|Y_i|],\qquad
 \lambda=\frac DP=Cd^3,\qquad
 A_0=\frac{3D+z\phi(z)}{2V},
\]
so that $g_i(u)\le\ell_i$ for all $u\in\R$ with equality on $\supp Y_i$.
Expanding around $u=0$ shows that uniformly in $i$ and $|u|\le H$,
\begin{equation}\label{eq:smooth-contact-bound}
       G_i(u)+b_iu+A_0u^2-\lambda|u|^3-\ell_i\le o(d)
\end{equation}
(since $C\to\ce>\phz/6$ makes the cubic coefficient $-\lambda+\phz d^3/6$
strictly negative). If $u\in\supp Y_i$ had $u+\E K_i$ bounded away from $\Z$,
then $\tau_i(u)$ would stay in $[\epsilon,1-\epsilon]$ (since $\{y+\mu\}\to0$)
and \eqref{eq:away-phase}--\eqref{eq:smooth-contact-bound} would give
$g_i(u)<\ell_i$, contradicting $u\in\supp Y_i$.
Thus $\eta_n=\max_i\|Z_i\|_\infty\to0$. Finally, $\E T^4\le3Q^2+\eta_n^2Q$
and $t=\E|T|\ge Q^{3/2}/(\E T^4)^{1/2}\ge Q/\sqrt{3Q+\eta_n^2}$ imply
$Q=o(t)$ whenever $t>0$.
\end{proof}

\section{Relative conditional first-moment bounds}\label{sec:relative}

To compare the global maximizer against its exact integer rounding $J_i$ when
$t=\E|T|>0$ is arbitrarily small, we first adjust the common scale by
$1+\delta=1+o(A^{-1/2})$ to place the full characteristic product at its exact
local maximum near $2\pi$, keeping the integer labels $K_i$ unchanged.

\begin{lemma}[Alignment at the exact full-product peak]\label{lem:full-peak}
After a common rescaling by $1+o(A^{-1/2})$, the summands $Y_i=J_i+Z_i$
retain the same integer parts $J_i$ and satisfy \eqref{eq:whole-phase} and
\begin{equation}\label{eq:relative-hypotheses}
\begin{gathered}
 A=\Var K\to\infty,\qquad |J_i|\le H=O(A^{1/3}),\\
 \sum_{i=1}^N\E|J_i|^3=O(A),\qquad Q\to0,\qquad
                 |\Cov(K,T)|\le KHQ.
\end{gathered}
\end{equation}
\end{lemma}
\begin{proof}
Because $|\E e^{\ii(2\pi+x/\sqrt A)\sum_i J_i}|\to e^{-x^2/2}$ and $T\to0$ in
$L^1$, the full product $|\E e^{\ii s\sum_i Y_i}|$ attains a local maximum at
$s=2\pi(1+\delta)$ with $\delta=o(A^{-1/2})$. Rescaling $Y_i$ by $1+\delta$
and realigning the phases $\beta_i=(2\pi)^{-1}\arg\E e^{2\pi\ii Z_i^{\rm new}}$
keeps the nearest-integer labels $K_i$ unchanged, while stationarity
$\sum_i\E[Y_i^{\rm new}\sin(2\pi E_i)]/r_i=0$ at the full peak yields
$|\Cov(\sum_i Y_i^{\rm new},T_{\rm new})|\le KHQ_{\rm new}$, proving
\eqref{eq:relative-hypotheses}.
\end{proof}

\begin{lemma}[Multiplier bound in terms of $t=\E|\sum_i Z_i|$]\label{lem:multiplier}
Let $Z_1,\ldots,Z_N$ be independent centered random variables with
$t=\E|\sum_{i=1}^N Z_i|$. For any random variables $W_1,\ldots,W_N$ such that
each pair $(W_i,Z_i)$ is independent of $(W_j,Z_j)_{j\ne i}$ and
$\|W_i\|_\infty\le M_W$,
\begin{equation}\label{eq:multiplier}
 \left|\sum_{i=1}^N\E[W_iZ_i]\right|
       \le8t\left(\sum_{i=1}^N\E W_i^2\right)^{1/2}+5M_Wt.
\end{equation}
\end{lemma}
\begin{proof}
Let $Z_i'$ be an independent copy of $Z_i$ (independent of $W_i$) and set
$D_i=Z_i-Z_i'$. By Khintchine's inequality, $\E(\sum_i D_i^2)^{1/2}\le2\sqrt3\,t=:C_0t$.
Truncating at $b=4C_0t$, Markov's inequality gives
$\Pp(\max_{j\ne i}|D_j|\le b)\ge3/4$, so summing over the disjoint events
where exactly one $|D_i|>b$ yields $\sum_i\E[|D_i|\ind_{\{|D_i|>b\}}]\le(4/3)C_0t<5t$.
For $U=\sum_i\E[D_i^2\ind_{\{|D_i|\le b\}}]$, second/fourth-moment
interpolation gives $U/\sqrt{U+b^2}\le C_0t$, so $U<60t^2$. Applying
$\E[W_iZ_i]=\E[W_iD_i]$, Cauchy--Schwarz on $\{|D_i|\le b\}$, and
$\|W_i\|_\infty\le M_W$ on $\{|D_i|>b\}$ proves \eqref{eq:multiplier}.
\end{proof}

\begin{lemma}[Relative conditional first-moment bounds]\label{lem:conditional-relative}
Under the hypotheses of Lemma~\ref{lem:full-peak}, whenever $t=\E|T|>0$, for
each fixed $B_0>0$ and uniformly over $|k-\mu|\le B_0\sqrt A$,
\begin{equation}\label{eq:conditional-relative}
 c\,t\le\E(|T|\mid K=k)\le Kt,\qquad
                      |\E(T\mid K=k)|=o(t).
\end{equation}
\end{lemma}
\begin{proof}
Take two independent copies $(K,T)$ and $(K',T')$ of the row, with
$D_i=Z_i-Z_i'$ and $D=\sum_i D_i$. Fix a large constant $B>1$, truncate at
$b=Bt$, and set $d_i=D_i\ind_{\{|D_i|\le b\}}$, $W=\sum_i d_i$,
$U=\sum_i\E d_i^2$, and $\Delta=\sum_i\E[|D_i|\ind_{\{|D_i|>b\}}]$. As in the
proof of Lemma~\ref{lem:multiplier},
\begin{equation}\label{eq:two-copy-truncation}
 \Delta\le Kt,\qquad \sum_{i=1}^N\Pp(|D_i|>b)\le \frac KB,\qquad
                      U\le K_Bt^2\quad\bigl(K_B=O(B)\bigr).
\end{equation}
On the frequency cell $\|s\|\le c/H$ ($s=(s_1,s_2)$), set
\[
 f_i(u)=\E e^{\ii uJ_i},\qquad
 m_{ir}(s)=\frac{\E[d_i^re^{\ii(s_1J_i+s_2J_i')}]}{f_i(s_1)f_i(s_2)}.
\]
Applying
Lemma~\ref{lem:multiplier} to the antisymmetric variables $d_i$ bounds the
linear twisted sum by
\begin{equation}\label{eq:relative-Fourier}
 \left|\sum_{i=1}^N m_{i1}(s)\right|\le Kt\Bigl\{H\|s\|
  +(\sqrt{HA}+H^2)\|s\|^2+(H^2\sqrt A+H^3)\|s\|^3\Bigr\},
\end{equation}
which is $o(t)$ on the Gaussian scale $\|s\|=O(A^{-1/2})$. Fourier inversion
of the second and fourth moments of $W$ (using \eqref{eq:torus-damping}
outside $\|s\|\le c/H$) yields, uniformly for $|k_1-\mu|,|k_2-\mu|\le B_0\sqrt A$,
\begin{equation}\label{eq:two-count-moments}
 \E(W^2\mid K=k_1,K'=k_2)=U+o(t^2),\qquad
 \E(W^4\mid K=k_1,K'=k_2)\le C_Bt^4.
\end{equation}
For the discarded tail $D-W=\sum_i D_i\ind_{\{|D_i|>b\}}$, one-pair deletion
ratios $\Pp(K_{-i}=k-k_{i1},K'_{-i}=k-k_{i2})/\Pp(K=K'=k)=1+o(1)$ uniformly
over shifts $O(H)$ imply $\E(|D-W|\mid K=K'=k)\le2\Delta\le2Kt$, while
\eqref{eq:two-count-moments} and $U\le K_Bt^2$ give
$\E(|W|\mid K=K'=k)\le K(B)t$, so $\E(|D|\mid K=K'=k)\le K'(B)t$. For the
lower bound, if $\Delta\ge\epsilon t$, conditioning on the disjoint events
where a single $|D_i|>b$ gives $\E(|D|\mid K=K'=k)\ge c\Delta\ge c\epsilon t$
for large $B$; if $\Delta<\epsilon t$, then $U\ge(1-\epsilon)^2t^2$ and
moment interpolation in \eqref{eq:two-count-moments} gives
$\E(|W|\mid K=K'=k)\ge c_Bt$, so choosing $\epsilon<c_B/4$ gives
$2\E(|T|\mid K=k)\ge\E(|D|\mid K=K'=k)\ge c_Bt/2$. Finally, applying the
one-copy Fourier expansion with \eqref{eq:relative-Fourier} and
$|\Cov(K,T)|\le KHQ=o(Ht)$ gives $|\E(T\mid K=k)|=o(t)$, and conditional
Jensen yields
\[
 \E(|T|\mid K=k)\le\E(|T-T'|\mid K=K'=k)+|\E(T\mid K=k)|\le K''(B)t.
\]
\end{proof}

\section{Exact phase collapse for global maximizers}\label{sec:comparison}

For the central integer layer $k=\operatorname{round}(y+\mu)$ and offset
$\theta=y+\mu-k\to0$, define
\[
 w(\theta)=\Pp(T>\theta\mid K=k),\qquad
                        f(\theta)=\theta+w(\theta).
\]

\begin{lemma}[Relative closed-threshold loss]\label{lem:closed-loss}
If $t=\E|T|>0$, then at the attaining central layer $k$ and offset $\theta$,
\begin{equation}\label{eq:relative-closed-loss}
        f(\theta)\ge ct,\qquad
                   |\theta|+w(\theta)+t\le K f(\theta).
\end{equation}
\end{lemma}
\begin{proof}
Truncating $T$ to $X=T\ind_{\{|T|\le1/4\}}$ changes its conditional CDF,
mean, and $L^1$ norm by $O(Q)=o(t)$ (via \eqref{eq:all-layer}), so
$m=\E(X\mid K=k)=o(t)$ and $s=\E(|X|\mid K=k)\ge ct$ by
Lemma~\ref{lem:conditional-relative}. Put $u=\Pp(X>\theta\mid K=k)$. For
$0\le\theta\le1/4$,
$s=2\E(X_+\mid K=k)-m\le2\theta+\frac12u-m\le2(\theta+u)-m$;
for $\theta=-b\in[-1/4,0)$, $m\le(1/4+b)u-b\le u/2-b$ and $s\le u/2-m$ give
$\theta+u\ge s+2m$, proving \eqref{eq:relative-closed-loss}.
\end{proof}

\begin{proposition}[Exact collapse of rounding residuals]\label{prop:collapse}
For every sufficiently late global maximizer in the sequence, $T=0$ almost
surely, so every summand $Y_i=J_i$ is supported on a single coset of $\Z$.
\end{proposition}
\begin{proof}
Suppose $t=\E|T|>0$ along a subsequence. Compare the global maximizer
$(Y_1,\ldots,Y_N)$ directly against the centered integer array
$(J_1,\ldots,J_N)$ at the same number of summands $N$ and integer threshold
$k-\mu$. By \eqref{eq:all-layer}, decomposing $\{S\le k+\theta-\mu\}$ by the
value of $K$ gives
\begin{equation}\label{eq:exact-crossing}
 F_Y(k+\theta-\mu)=F_K(k)-p_kw(\theta)+E_{\rm cross},\qquad
                         |E_{\rm cross}|\le KdQ=o\bigl(df(\theta)\bigr).
\end{equation}
Since $p_k=d\phi(z)+o(d)$, $\Phi((k+\theta-\mu)/\sqrt V)=\Phi((k-\mu)/\sqrt A)+d\phi(z)\theta+o(df(\theta))$
(using $|V-A|\le KHQ=o(At)$), and Lemma~\ref{lem:closed-loss} controls both
$|\theta|$ and $w(\theta)$ by $f(\theta)$, the discrepancies satisfy
\begin{equation}\label{eq:selected-comparison}
          D_Y=D_J-d\phi(z)f(\theta)+o\bigl(df(\theta)\bigr).
\end{equation}
Along the interpolation $Y_i(s)=J_i+sZ_i$ ($0\le s\le1$),
Lemma~\ref{lem:multiplier} with $W_i=3Y_i(s)|Y_i(s)|$ bounds
$\bigl|\frac{d}{ds}\sum_i\E|Y_i(s)|^3\bigr|\le K(\sqrt{HA}+H^2)t=o(At)$, so
$|P_Y-P_J|=o(At)$. Since $V\sim A$ ($d=V^{-1/2}$) and $P_Y\sim P_J\sim\re V$,
the prefactors $V^{3/2}/P_Y\sim(\re d)^{-1}$ and $A^{3/2}/P_J$ differ by a
relative error of $o(t)$, that is, $|V^{3/2}/P_Y-A^{3/2}/P_J|=o(d^{-1}t)$.
Because $D_J=O(d)$ and $t\le Kf(\theta)$, multiplying
\eqref{eq:selected-comparison} by $V^{3/2}/P_Y$ shows that
$(J_1,\ldots,J_N)$ has Berry--Esseen ratio strictly larger than that of
$(Y_1,\ldots,Y_N)$ by at least $c'f(\theta)\ge c''t>0$, contradicting global
maximality of $(Y_1,\ldots,Y_N)$. Thus $t=0$ eventually, so $Z_i=0$ almost
surely for all $i$.
\end{proof}

\section{Discrete concavity and reduction to two adjacent lattice points}\label{sec:lattice}

By Proposition~\ref{prop:collapse}, for all sufficiently late global
maximizers, each summand $Y_i$ is supported on a single coset of $\Z$.

\begin{lemma}[Reduction to two adjacent lattice points]\label{lem:adjacent}
For all sufficiently late global maximizers, every nonconstant summand $Y_i$
is supported on two adjacent points of its lattice coset.
\end{lemma}
\begin{proof}
Let $c_i(x)=\Pp(\sum_{j\ne i}Y_j=x)$ be the deleted point mass function on
its coset, with deleted variance $w_i=V-v_i$, $d_i=w_i^{-1/2}$, deleted
signed cubic sum $\kappa_i=\kappa(Y_{-i})$, and $H_3(x)=x^3-3x$. Fourier
inversion with multiplier $e^{-\ii u}-1$ and \eqref{eq:torus-damping} gives,
uniformly in $i$ and $x$,
\begin{equation}\label{eq:pmf-first-difference}
 c_i(x+1)-c_i(x)=d_i^2\phi'(d_ix)+\frac{d_i^3}2\phi''(d_ix)
       +d_i^3\frac{\kappa_i}{6w_i}(H_3\phi)'(d_ix)+o(d^3).
\end{equation}
For the contact function $g_i(u)=F_{-i}(y-u)+b_iu+A_0u^2-\lambda|u|^3$ on the
coset of $Y_i$, its centered second difference
$\Delta^2g_i(u)=g_i(u+1)-2g_i(u)+g_i(u-1)$ has deleted-CDF contribution
$c_i(y-u+1)-c_i(y-u)$, whose leading $-d^2z\phi(z)$ term cancels the
$d^2z\phi(z)$ part of $2A_0$, leaving for $|u|\le H+1$:
\begin{equation}\label{eq:lattice-concavity}
 d^{-3}\Delta^2g_i(u)=\phz(1+u)-\ce
       \bigl\{|u+1|^3+|u-1|^3-2|u|^3\bigr\}
       +o(1+|u|)+O(d|z|v_i).
\end{equation}
Because $\ce>\phz$ and $|u+1|^3+|u-1|^3-2|u|^3$ equals $6|u|$ for $|u|\ge1$
and $2+6u^2-2|u|^3\ge2+4u^2$ for $|u|\le1$, the leading expression in
\eqref{eq:lattice-concavity} is bounded above by $-c(1+|u|)$ with $c>0$.
Applying the discrete Green identity on any interval of length $L\ge2$
between two contact points $u_0,u_0+L\in\supp Y_i$,
\[
 g_i(u_0+j)-g_i(u_0)
 =-\sum_{k=1}^{L-1}\frac{\min(j,k)(L-\max(j,k))}{L}\Delta^2g_i(u_0+k)\le0,
\]
first bounds $L$ and $v_i$ by a uniform constant (making $d|z|v_i=o(1)$ in
\eqref{eq:lattice-concavity}), and then implies $\Delta^2g_i(u)<0$ strictly
at every coset point, which makes $L\ge2$ impossible. Thus $L=1$ for every
nonconstant summand.
\end{proof}

Since every nonconstant summand is supported on two points at distance $1$,
all active summands share a common support diameter. Applying Theorem~1.1 of
\citet{BellBernoulli2026} at the same horizon $N_n$ shows that all active
summands are identically distributed, completing the proof of
Theorem~\ref{thm:global}.


\section{A large binary group}\label{sec:large-group-example}
As an example of Theorem~\ref{thm:global}, suppose an actual full global
array at horizon $M+r_M$ contains $M\to\infty$ active binary laws of a
common diameter $h_M$, with probabilities in a fixed compact
$J\subset(0,1)$. Let
\[
 R_{\rm rem}=\sum_{i\text{ outside the group}}\E|X_i|^3.
\]
The direct original-price condition
\begin{equation}\label{eq:large-group-price}
                      R_{\rm rem}=o(M^{3/2}h_M^3)
\end{equation}
is sufficient. The given group contributes variance at least
$c_JMh_M^2$ and price at most $KMh_M^3$, so
\[
 L\le\frac{KMh_M^3+R_{\rm rem}}
                 {(c_JMh_M^2)^{3/2}}\longrightarrow0.
\]
Theorem~\ref{thm:global} therefore makes every sufficiently late original
array padded IID binary. Its common diameter is $h_M$, because the
original group is active: each remaining coordinate is either the
same binary law or zero. No bound on the number of remaining
coordinates is needed beyond their stated additive price.

For instance, if all remaining centered laws have support diameter at
most $H_Mh_M$, then
$R_{\rm rem}\le r_MH_M^3h_M^3$.
Thus $r_MH_M^3=o(M^{3/2})$ implies
\eqref{eq:large-group-price} directly.

\begingroup\small\raggedright
\bibliographystyle{plainnat}\bibliography{refs}
\endgroup
\end{document}
